\textbf{Proof.}
Put \(m=m_1\wedge m_0\).  We first establish the rank-uniform matrix bound.  For every \([B]_+\in\mathcal B_Z^+\), lem:rank-stratified-projector-reduction supplies the orthogonal projector
\[
P_B=\Omega^{1/2}B(B^{\top}\Omega B)^+B^{\top}\Omega^{1/2}
\]
onto \(\operatorname{range}(\Omega^{1/2}B)\).  Thus \(0\preceq P_B\preceq I_{2n}\) for ranks zero, one, and two.  Congruence by \(\Pi^{-1}\Omega^{1/2}\) gives
\[
0\preceq
A_B^+=\Pi^{-1}\Omega^{1/2}P_B\Omega^{1/2}\Pi^{-1}
\preceq
\Pi^{-1}\Omega\Pi^{-1}=A_{\mathrm{HT}}.
\tag{1}
\]
Both matrices in (1) are positive semidefinite, so their nuclear norms equal their traces.  Since \(D_{i,a}\) is Bernoulli with mean \(\pi_{i,a}\),
\[
\Omega_{(i,a),(i,a)}
=\operatorname{Var}_{\nu}(D_{i,a})
=\pi_{i,a}(1-\pi_{i,a}).
\]
Consequently, with \(\pi_{\min,n}=\min_{i,a}\pi_{i,a}>0\) by ass:exposure-positivity,
\[
\begin{aligned}
\lVert A_B^+\rVert_*
&\le \operatorname{tr}(A_{\mathrm{HT}})
=\sum_{i=1}^n\sum_{a\in\{1,0\}}
\frac{1-\pi_{i,a}}{\pi_{i,a}}
\le 2n\pi_{\min,n}^{-1},\\
\lVert A_{\mathrm{HT}}\rVert_*
&\le 2n\pi_{\min,n}^{-1}.
\end{aligned}
\tag{2}
\]

For \(e\in\{L,U\}\), def:pilot-plugin and def:completion-pseudoinverse-loss give
\[
\widehat\ell_{B,e}^+-\ell_B^+(c_e)
=n^{-2}\operatorname{tr}\!\left[
(A_{\mathrm{HT}}-A_B^+)\{\widehat Q_e-Q(c_e)\}
\right].
\]
Trace duality and (2) imply
\[
\left|\widehat\ell_{B,e}^+-\ell_B^+(c_e)\right|
\le 4\pi_{\min,n}^{-1}n^{-1}M_m.
\tag{3}
\]
The map \((x_L,x_U)\mapsto\max\{x_L,x_U\}\) is one-Lipschitz for the sup norm.  Taking the maximum over endpoints and then the supremum over all three rank strata in (3) proves
\[
d_m^+\le 4\pi_{\min,n}^{-1}n^{-1}M_m.
\tag{4}
\]
Along a triangular sequence satisfying ass:zhao-exposure-probability-rate, \(\pi_{\min,n}^{-1}\le\underline\pi^{-1}n^{\beta}\), so (4) is exactly (A).

We next derive directly, without a basis-class premise, the endpoint-only pilot rate needed below.  For an arm \(a\), let
\[
D_a=W_1(\widehat F_a,F_a)
=\int_0^1|\widehat F_a(y)-F_a(y)|\,dy
=\int_0^1|\widehat F_a^{-1}(u)-F_a^{-1}(u)|\,du.
\tag{5}
\]
The equalities in (5) follow from the layer-cake identity and Fubini's theorem.  At each \(y\), the empirical distribution function is an average of \(m_a\) Bernoulli variables.  Jensen's inequality and their variance formula therefore yield
\[
\mathbb E_P D_a
\le\int_0^1
\sqrt{\frac{F_a(y)\{1-F_a(y)\}}{m_a}}\,dy
\le\frac{1}{2\sqrt{m_a}}.
\tag{6}
\]
Because \(y\mapsto y\) and \(y\mapsto y^2\) are respectively one- and two-Lipschitz on \([0,1]\),
\[
|\widehat\mu_a-\mu_a|\le D_a,
\qquad
|\widehat\sigma_a^2-\sigma_a^2|\le4D_a,
\qquad
|\widehat\mu_a^2-\mu_a^2|\le2D_a.
\tag{7}
\]
For either endpoint coupling, the inequality \(|xy-x'y'|\le|x-x'|+|y-y'|\) on \([0,1]^4\), applied to the comonotone or countermonotone quantiles, gives
\[
|\widehat c_e-c_e|\le2(D_1+D_0),
\qquad e\in\{L,U\}.
\tag{8}
\]
Indeed, one copy of \(D_1+D_0\) bounds the endpoint cross-moment error and another bounds the product-of-means error.

Using the block formula in def:moment-matrix, \(\lVert I_n\rVert_{\mathrm{op}}=1\), and \(\lVert\mathbf 1\mathbf 1^{\top}\rVert_{\mathrm{op}}=n\), equations (7)--(8) imply
\[
\lVert\widehat Q_e-Q(c_e)\rVert_{\mathrm{op}}
\le (4n+8)(D_1+D_0)
\le 7n(D_1+D_0),
\tag{9}
\]
where the last inequality uses \(n\ge3\).  Combining (6) and (9), there is the numerical constant \(K_M=7\) such that
\[
\mathbb E_P M_m\le K_M\frac{n}{\sqrt m}.
\tag{10}
\]
This derivation uses only the bounded unpaired marginal samples and is independent of every basis class and every \(\kappa\)-feasibility condition.

By the defining positive-tolerance property of \(\widehat B_{m,\varepsilon_m}^+\), and by the definition of \(d_m^+\), pointwise in the pilot sample,
\[
\begin{aligned}
0\le w^+(\widehat B_{m,\varepsilon_m}^+)-v_{\mathrm{mm}}^+
&\le \widehat w_m^+(\widehat B_{m,\varepsilon_m}^+)
   +d_m^+-v_{\mathrm{mm}}^+\\
&\le \inf_{[B]_+\in\mathcal B_Z^+}\widehat w_m^+(B)
   +\varepsilon_m+d_m^+-v_{\mathrm{mm}}^+\\
&\le 2d_m^++\varepsilon_m.
\end{aligned}
\tag{11}
\]
The last inequality follows from
\[
\left|
\inf_{[B]_+\in\mathcal B_Z^+}\widehat w_m^+(B)
-\inf_{[B]_+\in\mathcal B_Z^+}w^+(B)
\right|\le d_m^+.
\]
Taking pilot expectations in (11) and applying (A) proves (B).  Finally, (10) gives
\[
\begin{aligned}
n\mathbb E_P\!\left[
w^+(\widehat B_{m,\varepsilon_m}^+)-v_{\mathrm{mm}}^+
\right]
&\le
8K_M\underline\pi^{-1}
\frac{n^{1+\beta}}{\sqrt{m_1\wedge m_0}}
+n\mathbb E_P[\varepsilon_m]
\longrightarrow0
\end{aligned}
\]
under the two displayed hypotheses.  Every matrix bound covered ranks zero, one, and two.  No step constructs or analyzes an ordinary-inverse observable estimator for a singular selected basis, so the conclusion is solely a full-carrier basis-objective regret result. \(\square\)